% Part: turing-machines
% Chapter: undecidability
% Section: verification

\documentclass[../../../include/open-logic-section]{subfiles}

\begin{document}

\olfileid[fr]{tur}{und}{ver}
\olsection{Vérifier la représentation}

\begin{explain}
Pour vérifier que notre représentation fonctionne, nous devons
établir deux implications. Si~$M$ s'arrête sur l'entrée~$w$, alors
$!T(M,w)\lif !E(M,w)$ est valide. Réciproquement, si cet énoncé est
valide, alors~$M$ finit effectivement par s'arrêter sur~$w$.

Les stratégies diffèrent. Pour la première implication, il faut
montrer qu'!!a{sentence} du premier ordre, à savoir
$!T(M,w)\lif !E(M,w)$, est valide. On peut établir cette validité
en donnant !!a{derivation}. Mais l'argument doit valoir pour tous
$M$ et~$w$ : il porte donc sur une infinité d'!!{sentence}s, et
non sur un seul dont nous pourrions exhiber !!a{derivation} fixe.
Nous allons montrer par récurrence que, quels que soient~$M$ et~$w$
et quel que soit le nombre d'étapes avant l'arrêt, il existe
!!a{derivation} de $!T(M,w)\lif !E(M,w)$.

La récurrence suit le nombre d'étapes parcourues avant d'atteindre
une configuration d'arrêt. Nous montrerons que, pour chaque
instant~$n$ atteint dans l'exécution de~$M$ sur~$w$,
$!T(M,w)\Entails !C(M,w,n)$, où $!C(M,w,n)$ décrit correctement
la configuration après~$n$ étapes. Si~$M$ s'arrête après~$n$
étapes, cette configuration est terminale. Nous établirons aussi
que $!C(M,w,n)\Entails !E(M,w)$ pour une telle configuration.
Ainsi, si~$M$ s'arrête, il existe un~$n$ tel que
$!T(M,w)\Entails !C(M,w,n)\Entails !E(M,w)$; autrement dit,
$\Entails !T(M,w)\lif !E(M,w)$. La complétude fournit également
une dérivation de cette implication.

Pour la réciproque, supposons $\Entails !T(M,w)\lif !E(M,w)$.
Cela signifie que $!T(M,w)\Entails !E(M,w)$ : $!E(M,w)$ est vrai
dans toute !!{structure} où $!T(M,w)$ est vrai. Nous construirons
une !!{structure}~$\Struct M$ de domaine~$\Nat$ satisfaisant
$!T(M,w)$, dans laquelle les interprétations de $\Obj Q_q$ et
de~$\Obj S_\sigma$ sont déterminées par l'exécution réelle de~$M$
sur~$w$. Par exemple,
$\Sat{M}{\Obj Q_q(\num{m},\num{n})}$ signifie qu'après $n$ étapes,
$M$ est dans l'état~$q$ et sa tête lit la case~$m$. Puisque
$\Sat{M}{!T(M,w)}$ par construction et que
$!T(M,w)\Entails !E(M,w)$, nous aurons $\Sat{M}{!E(M,w)}$.
Il existera alors un instant~$n\in\Domain{M}=\Nat$ auquel~$M$
se trouve dans une configuration d'arrêt.
\end{explain}

\begin{defn}
Soit~$!C(M,w,n)$ l'!!{sentence}
\[
\Obj Q_q(\num{m},\num{n})\land
\Obj S_{\sigma_0}(\num{0},\num{n})\land\dots\land
\Obj S_{\sigma_k}(\num{k},\num{n})\land
\lforall[x][(\num{k}<x\lif\Obj S_\TMblank(x,\num{n}))]
\]
où $q$ est l'état de~$M$ à l'instant~$n$, sa tête lit alors la
case~$m$, la case~$i$ contient~$\sigma_i$ pour $0\le i\le k$,
et $k$ est le plus grand des deux indices suivants :
$\len{w}$, indice de la dernière case explicitement spécifiée
par l'entrée initiale (ou du marqueur si $w$ est vide), et
l'indice de la case la plus à droite que la tête a visitée
jusqu'à l'instant~$n$.
\end{defn}

\begin{lem}\ollabel{lem:halt-config-implies-halt}
Si l'exécution de~$M$ sur~$w$ atteint une configuration d'arrêt
après $n$ étapes, alors $!C(M,w,n)\Entails !E(M,w)$.
\end{lem}

\begin{proof}
Supposons que~$M$ s'arrête après $n$ étapes. Il existe un état~$q$,
une case~$m$ et un symbole~$\sigma$ tels que :
\begin{enumerate}
\item après $n$ étapes, $M$ est dans l'état~$q$ et lit la case~$m$,
  qui contient~$\sigma$;
\item la fonction de transition $\delta(q,\sigma)$ n'est pas définie.
\end{enumerate}
La description~$!C(M,w,n)$ contient donc les conjoints
$\Obj Q_q(\num{m},\num{n})$ et
$\Obj S_\sigma(\num{m},\num{n})$. Leur conjonction entraîne
$!E(M,w)$ :
\[
\lexists[x][\lexists[y][(\bigvee_{\tuple{q,\sigma}\in X}
  (\Obj Q_q(x,y)\land\Obj S_\sigma(x,y)))]]
\]
En effet, le couple état-symbole réellement lu appartient à~$X$,
et $\Obj Q_{q'}(\num{m},\num{n})\land
\Obj S_{\sigma'}(\num{m},\num{n})$ entraîne le terme correspondant
de la disjonction lorsque $\tuple{q',\sigma'}\in X$.
\end{proof}

\begin{explain}
Si~$M$ s'arrête sur~$w$, il existe donc un~$n$ tel que
$!C(M,w,n)\Entails !E(M,w)$. Montrons maintenant que
$!T(M,w)\Entails !C(M,w,n)$ pour chaque instant~$n$ atteint
dans l'exécution.
\end{explain}

\begin{lem}
\ollabel{lem:config}
Pour chaque~$n$ tel que l'exécution de~$M$ sur~$w$ atteint une
configuration après $n$ étapes, $!T(M,w)\Entails !C(M,w,n)$.
La configuration après~$n$ étapes peut être une configuration d'arrêt.\footnote{Corrections éditoriales : la source exclut cet instant de l'énoncé du lemme, alors que la suite de la preuve l'utilise. La borne $k$ de la configuration comprend toute l'entrée explicitement donnée, même si elle finit par des blancs ; la preuve traite aussi l'entrée vide et le premier pas.}
\end{lem}

\begin{proof}
Procédons par récurrence sur~$n$. Pour $n=0$, l'état initial et
les symboles des cases~$0$ à~$\len{w}$ sont donnés par les
axiomes de~$!T(M,w)$; les cases au-delà sont blanches. Si
l'entrée est vide, la tête lit néanmoins la case~$1$ : sa
blancheur découle de l'axiome de blancheur et de
$\num{0}<\num{1}$. On obtient ainsi tous les conjoints de
$!C(M,w,0)$, y compris l'énoncé de blancheur à droite de la
borne~$k$ choisie dans sa définition.

Supposons que $!T(M,w)\Entails !C(M,w,n)$ et que la configuration
après $n$ étapes ne soit pas terminale. Nous devons montrer que
$!T(M,w)\Entails !C(M,w,n+1)$. Soit~$q$ l'état à l'instant~$n$,
$m$ la case lue et~$\sigma$ son symbole. La fonction de transition
est définie en~$\tuple{q,\sigma}$ et donne une instruction de
l'une des trois formes suivantes :
\begin{enumerate}
\item \ollabel{right}
  $\delta(q,\sigma)=\tuple{q',\sigma',\TMright}$;
\item \ollabel{left}
  $\delta(q,\sigma)=\tuple{q',\sigma',\TMleft}$;
\item \ollabel{stay}
  $\delta(q,\sigma)=\tuple{q',\sigma',\TMstay}$.
\end{enumerate}
Examinons ces cas.

\begin{enumerate}
\item Supposons l'instruction~\olref{right}. D'après
  la \olref[rep]{defn:tm-descr}\olref[rep]{rep-right}, $!T(M,w)$
  contient l'énoncé
\begin{align*}
& \lforall[x][\lforall[y][(
  (\Obj Q_q(x,y)\land\Obj S_\sigma(x,y))\lif {}]]\\
&\qquad (\Obj Q_{q'}(x',y')\land
  \Obj S_{\sigma'}(x,y')\land !A(x,y))).
\end{align*}
  Par instanciation universelle, en remplaçant $x$ par~$\num{m}$
  et $y$ par~$\num{n}$, nous obtenons l'!!{sentence}
\begin{align*}
& (\Obj Q_q(\num{m},\num{n})\land
  \Obj S_\sigma(\num{m},\num{n}))\lif {}\\
&\qquad (\Obj Q_{q'}(\num{m}',\num{n}')\land
  \Obj S_{\sigma'}(\num{m},\num{n}')\land
  !A(\num{m},\num{n})).
\end{align*}
  Par l'hypothèse de récurrence, $!T(M,w)$ entraîne
  $!C(M,w,n)$, donc en particulier les deux conjoints de
  l'antécédent. Par modus ponens, il entraîne
  $\Obj Q_{q'}(\num{m}',\num{n}')$,
  $\Obj S_{\sigma'}(\num{m},\num{n}')$ et
  $!A(\num{m},\num{n})$.

  Pour chaque $i\le k$ distinct de~$m$, les axiomes d'ordre
  établissent $\num{i}<\num{m}$ ou
  $\num{m}<\num{i}$; la formule~$!A$ conserve donc le symbole
  de la case~$i$. Nous obtenons les faits sur les cases
  \emph{autres que celle où l'instruction écrit}\footnote{Correction éditoriale : l'ellipse du calcul source inclut encore le symbole ancien de la case $m$, bien que le commentaire l'exclue. Nous écrivons la conjonction avec $i\ne m$ et traitons séparément le symbole nouveau.} :
\[
\bigwedge_{\substack{0\le i\le k\\i\ne m}}
  \Obj S_{\sigma_i}(\num{i},\num{n}').
\]
  Le symbole de la case~$m$ est celui, déjà établi,
  $\sigma'$. De même, $!A(\num{m},\num{n})$ conserve la
  blancheur à droite de~$k$. En effet, si $m<k$,
  la \olref[rep]{prop:mlessk} et la transitivité de~$<$ donnent
  $\lforall[x][(\num{k}<x\lif\num{m}<x)]$; si $m=k$,
  cette implication est immédiate. En la combinant avec
  $\lforall[x][(\num{k}<x\lif
  \Obj S_\TMblank(x,\num{n}))]$ et avec~$!A$, nous obtenons
  $\lforall[x][(\num{k}<x\lif
  \Obj S_\TMblank(x,\num{n}'))]$.

  Si $m<k$, ces faits décrivent déjà $!C(M,w,n+1)$.
  Si $m=k$, la tête a aussi visité la case~$k+1$ après la
  nouvelle étape : la borne de la description devient~$k+1$.
  Il faut établir deux !!{sentence}s supplémentaires :
  $\Obj S_\TMblank(\num{k}',\num{n}')$
  et $\lforall[x][(\num{k}'<x\lif
  \Obj S_\TMblank(x,\num{n}'))]$. Le premier fait découle de
  $\num{k}<\num{k}'$ et de la blancheur déjà obtenue au-delà
  de~$k$. Le second suit de la transitivité de~$<$; la dernière
  étape de ce raisonnement est proposée en exercice ci-dessous.

\item Supposons l'instruction~\olref{left}. D'après
  la \olref[rep]{defn:tm-descr}\olref[rep]{rep-left}, $!T(M,w)$
  contient la conjonction suivante, avec le terme de conservation
  $!A(x',y)$ corrigé pour exclure la case où l'on écrit :
\begin{align*}
& \lforall[x][\lforall[y][(
  (\Obj Q_q(x',y)\land\Obj S_\sigma(x',y))\lif {}]]\\
&\qquad (\Obj Q_{q'}(x,y')\land
  \Obj S_{\sigma'}(x',y')\land !A(x',y)))\land {}\\
& \lforall[y][(
  (\Obj Q_q(\num{0},y)\land
  \Obj S_\sigma(\num{0},y))\lif {}]\\
&\qquad (\Obj Q_{q'}(\num{0},y')\land
  \Obj S_{\sigma'}(\num{0},y')\land
  !A(\num{0},y))).
\end{align*}
  Si $m>0$, posons $l=m-1$, de sorte que $m=l+1$.
  Le premier conjoint de cet !!{sentence} donne, par instanciation,
\begin{align*}
& (\Obj Q_q(\num{l}',\num{n})\land
  \Obj S_\sigma(\num{l}',\num{n}))\lif {}\\
&\qquad (\Obj Q_{q'}(\num{l},\num{n}')\land
  \Obj S_{\sigma'}(\num{l}',\num{n}')\land
  !A(\num{l}',\num{n})).
\end{align*}
  Si $m=0$, posons $l=0$. Le second conjoint donne l'!!{sentence}
\begin{align*}
& (\Obj Q_q(\num{0},\num{n})\land
  \Obj S_\sigma(\num{0},\num{n}))\lif {}\\
&\qquad (\Obj Q_{q'}(\num{0},\num{n}')\land
  \Obj S_{\sigma'}(\num{0},\num{n}')\land
  !A(\num{0},\num{n})).
\end{align*}
  Dans les deux cas, l'hypothèse de récurrence fournit
  l'antécédent de l'implication applicable. Nous obtenons donc
  $\Obj Q_{q'}(\num{l},\num{n}')$ et
  $\Obj S_{\sigma'}(\num{m},\num{n}')$.
  Comme $\num{l}'\ident\num{m}$ lorsque $m>0$, la formule
  de conservation est alors $!A(\num{m},\num{n})$ dans les deux
  cas. Elle transporte tous les symboles des autres cases et la
  blancheur au-delà de~$k$, comme pour le déplacement à droite.
  La tête s'étant déplacée vers la gauche ou étant restée en~$0$,
  la borne~$k$ ne change pas. L'ensemble de ces faits est
  précisément $!C(M,w,n+1)$.

\item Le cas de l'instruction~\olref{stay} est laissé en exercice.
\end{enumerate}
La récurrence établit donc
$!T(M,w)\Entails !C(M,w,n)$ à chaque instant atteint par
l'exécution, y compris celui de l'arrêt éventuel.
\end{proof}

\begin{prob}
Complétez le cas~\olref[tur][und][ver]{stay} dans la preuve du
\olref[tur][und][ver]{lem:config}.
\end{prob}

\begin{prob}
Donnez !!a{derivation} de
$\Obj S_{\sigma_i}(\num{i},\num{n}')$ à partir de
$\Obj S_{\sigma_i}(\num{i},\num{n})$ et de~$!A(\num{m},\num{n})$,
en supposant $i\ne m$, c'est-à-dire $i<m$ ou $m<i$.
\end{prob}

\begin{prob}
Donnez !!a{derivation} de
$\lforall[x][(\num{k}'<x\lif
\Obj S_\TMblank(x,\num{n}'))]$ à partir de
$\lforall[x][(\num{k}<x\lif
\Obj S_\TMblank(x,\num{n}'))]$,
de $\lforall[x][x<x']$ et de
$\lforall[x][\lforall[y][\lforall[z][
((x<y\land y<z)\lif x<z)]]]$.
\end{prob}

\begin{lem}
\ollabel{lem:valid-if-halt}
Si~$M$ s'arrête sur l'entrée~$w$, alors
$!T(M,w)\lif !E(M,w)$ est valide.
\end{lem}

\begin{proof}
D'après le \olref{lem:config}, la description~$!C(M,w,n)$ de
chaque configuration atteinte après $n$ étapes découle de~$!T(M,w)$.
Supposons que~$M$ s'arrête après $k$ étapes. Sa tête lit alors
une case~$m\in\Nat$, dans un état~$q$, et cette case contient
un symbole~$\sigma$ tel que $\delta(q,\sigma)$ ne soit pas
définie. Ainsi, $!C(M,w,k)$ décrit une configuration d'arrêt
et contient $\Obj Q_q(\num{m},\num{k})$ ainsi que
$\Obj S_\sigma(\num{m},\num{k})$. D'après le
\olref{lem:halt-config-implies-halt},
$!C(M,w,k)\Entails !E(M,w)$. Puisque
$!T(M,w)\Entails !C(M,w,k)$, il s'ensuit
$!T(M,w)\Entails !E(M,w)$, donc la validité de
$!T(M,w)\lif !E(M,w)$.
\end{proof}

\begin{explain}
Il reste à prouver la réciproque : si
$!T(M,w)\lif !E(M,w)$ est valide, alors~$M$ s'arrête
effectivement sur l'entrée~$w$.
\end{explain}

\begin{lem}
\ollabel{lem:halt-if-valid}
Si $\Entails !T(M,w)\lif !E(M,w)$, alors~$M$ s'arrête
sur l'entrée~$w$.
\end{lem}

\begin{proof}
Considérons la $\Lang L_M$-!!{structure}~$\Struct M$ de
domaine~$\Nat$ qui interprète $\Obj 0$ par~$0$,
$\prime$ par la fonction successeur et~$<$ par l'ordre usuel.
Nous interprétons les prédicats $\Obj Q_q$ et~$\Obj S_\sigma$
au moyen de l'exécution effective de~$M$ sur~$w$ :
\begin{align*}
\Assign{\Obj Q_q}{M}
&=\Setabs{\tuple{m,n}}{\begin{array}{l}
  \text{après $n$ étapes sur $w$, $M$ est dans l'état $q$}\\
  \text{et sa tête lit la case $m$}
  \end{array}},\\
\Assign{\Obj S_\sigma}{M}
&=\Setabs{\tuple{m,n}}{\begin{array}{l}
  \text{après $n$ étapes sur $w$, la case $m$}\\
  \text{contient le symbole $\sigma$}
  \end{array}}.
\end{align*}
Si l'exécution s'arrête, les relations ainsi définies ne
contiennent aucun couple correspondant à un instant ultérieur.
Les axiomes de transition de~$!T(M,w)$ ne prescrivent alors
aucune étape supplémentaire : leurs antécédents y sont faux.

La !!{structure}~$\Struct M$ décrit donc l'exécution réelle,
étape par étape, et $\Sat{M}{!T(M,w)}$. Si
$\Entails !T(M,w)\lif !E(M,w)$, alors aussi
$\Sat{M}{!E(M,w)}$, c'est-à-dire
\[
\Sat{M}{\lexists[x][\lexists[y][(\bigvee_{\tuple{q,\sigma}\in X}
  (\Obj Q_q(x,y)\land\Obj S_\sigma(x,y)))]]}.
\]
Puisque $\Domain{M}=\Nat$, il existe donc $m,n\in\Nat$ et
un couple~$\tuple{q,\sigma}\in X$ tels que
$\Sat{M}{\Obj Q_q(\num{m},\num{n})\land
\Obj S_\sigma(\num{m},\num{n})}$. Par la définition de
$\Struct M$, cela signifie qu'après $n$ étapes sur~$w$, la
machine~$M$ est dans l'état~$q$ et lit le symbole~$\sigma$.
Or $\tuple{q,\sigma}\in X$ signifie que
$\delta(q,\sigma)$ n'est pas définie. La machine s'est
donc arrêtée.
\end{proof}

\end{document}
